\begin{proof}[Proof of \cref{obj:thm:uniform-frontier}]
Fix public constants \(K=(p_{\min},c_{\mathrm{lo}},c_{\mathrm{hi}})\), \(\beta\in(0,1]\), \(\kappa\in[0,2]\), and \(\alpha\in(0,1/2)\), under the envelope calibration in the statement. Throughout, use the rates and comparison scales defined there. In particular,
\[
d=2\beta+\kappa+1>1>0.
\]

\begin{enumerate}
\item \textbf{Upper risk bound.}
By \cref{obj:thm:observable-upper}, there are \(C_R>0\) and \(n_R\in\mathbb N\), independent of the public path space and of \(\sigma\), such that
\[
\sup_{P\in\mathcal P_{K,\beta,\kappa,\sigma}}
E_{Q_P^n}|\widehat\theta-\theta(P)|
\le C_R\rho_{n,\sigma}
\qquad
(n\ge n_R,\ 0\le\sigma\le1/4).
\]
The measurable estimator in \cref{obj:def:total-estimator} is admissible in \cref{obj:def:minimax-risk}. Viewed as a randomized estimator, it ignores the independent uniform seed, so its risk under the product experiment equals its observed-sample risk. Consequently,
\[
R_n(\sigma)
\le
\sup_{P\in\mathcal P_{K,\beta,\kappa,\sigma}}
E_{Q_P^n}|\widehat\theta-\theta(P)|
\le C_R\rho_{n,\sigma}.
\]
% lean: CausalSmith.Stat.NoisydoseWeakdesignTransition.minimaxRisk_frontier_upper

\item \textbf{Upper length bound.}
By \cref{obj:thm:finite-honesty}, the measurable connected interval in \cref{obj:def:honest-interval} satisfies
\[
Q_P^n\{\theta(P)\in C_{n,\sigma}\}\ge1-\alpha
\]
for every model law, every \(n\), and every allowed \(\sigma\). The same result supplies \(C_H>0\) and \(n_H\in\mathbb N\), uniform over the path spaces and error scales, such that
\[
\sup_{P\in\mathcal P_{K,\beta,\kappa,\sigma}}
E_{Q_P^n}\operatorname{length}(C_{n,\sigma})
\le C_H\rho_{n,\sigma}
\qquad(n\ge n_H).
\]
Thus this procedure belongs to the collection defining \cref{obj:def:honest-length}, and
\[
H_n(\sigma)
\le
\sup_{P\in\mathcal P_{K,\beta,\kappa,\sigma}}
E_{Q_P^n}\operatorname{length}(C_{n,\sigma})
\le C_H\rho_{n,\sigma}.
\]
% lean: CausalSmith.Stat.NoisydoseWeakdesignTransition.honestLength_frontier_upper

\item \textbf{A common lower bound.}
Use the published Jacobi localization bound
\citep[Theorem 2.4, equation (2.14)]{PetrushevXu2005},
the local kernel Lipschitz bound
\citep[Theorem 2.2, equation (2.4)]{https://people.math.sc.edu/pencho/Publications/kpx-06-28-06-web.pdf},
Jacobi orthogonality and the positive squared-norm formula
\citep[Table 18.3.1, Jacobi row and caption]{https://dlmf.nist.gov/18.3},
and the positive-real-axis gamma-ratio asymptotic
\citep[equation (5.11.12)]{https://dlmf.nist.gov/5.11.E12}
as the analytic inputs to \cref{obj:thm:observed-lower}.
The packet parameters are admissible: for \(\kappa>0\), they are
\(r=4\) and \(b=(\kappa-1)/2>-1/2\), while the \(\kappa=0\) construction uses the symmetric parameters \((4,4)\).

Set
\[
\tau:=\min\left\{\frac15,\frac{1-2\alpha}{2}\right\}.
\]
Since \(\alpha<1/2\), this choice gives
\[
0<\tau,\qquad \tau\le\frac15,\qquad
\tau\le1-2\alpha-\tau.
\]
That result supplies \(c_{\mathrm w}>0\) and \(n_{\mathrm w}\in\mathbb N\), uniform over the path spaces, such that for every \(n\ge n_{\mathrm w}\) and every allowed \(\sigma\), there are model laws \(P_+,P_-\) with
\[
\theta_\pm:=\theta(P_\pm),\qquad
\delta_{n,\sigma}:=c_{\mathrm w}\rho_{n,\sigma},
\]
\[
|\theta_+-\theta_-|\ge\delta_{n,\sigma},
\qquad
\operatorname{TV}(Q_{P_+}^n,Q_{P_-}^n)\le\tau.
\]
For \(n\ge1\), \(L_n\ge1\); each branch defining \(b_{n,\sigma}\) is nonnegative when \(\sigma\ge0\). Hence
\[
\rho_{n,\sigma}\ge0,\qquad \delta_{n,\sigma}\ge0.
\]

To handle the permitted randomization, define the seed law and product laws by
\[
\nu_{\mathrm{seed}}:=\operatorname{Unif}[0,1],
\qquad
Q_\pm^{n,\mathrm{rand}}:=Q_{P_\pm}^n\otimes\nu_{\mathrm{seed}}.
\]
For any measurable event \(\mathcal V\) in the sample--seed product space, define its section probability by
\[
f_{\mathcal V}(z)
:=
\nu_{\mathrm{seed}}\{u_{\mathrm{seed}}\in[0,1]:
(z,u_{\mathrm{seed}})\in\mathcal V\}.
\]
This is measurable and lies in \([0,1]\). Product integration and the layer-cake identity therefore give
\[
\begin{aligned}
\bigl|Q_+^{n,\mathrm{rand}}(\mathcal V)
      -Q_-^{n,\mathrm{rand}}(\mathcal V)\bigr|
&=
\left|\int f_{\mathcal V}\,dQ_{P_+}^n
      -\int f_{\mathcal V}\,dQ_{P_-}^n\right|\\
&=
\left|\int_0^1
\left[
Q_{P_+}^n\{f_{\mathcal V}\ge v\}
-
Q_{P_-}^n\{f_{\mathcal V}\ge v\}
\right]\,dv\right|\\
&\le \operatorname{TV}(Q_{P_+}^n,Q_{P_-}^n).
\end{aligned}
\]
Taking the supremum over such events yields
\[
\operatorname{TV}(Q_+^{n,\mathrm{rand}},Q_-^{n,\mathrm{rand}})
\le\tau\le\frac15.
\]

For any measurable real-valued estimator \(T\) on this product experiment, define
\[
\mathcal A_\pm
:=
\left\{(z,u_{\mathrm{seed}}):
|T(z,u_{\mathrm{seed}})-\theta_\pm|
\ge\frac{\delta_{n,\sigma}}2\right\}.
\]
The separation and triangle inequality imply
\(\mathcal A_+^c\subseteq\mathcal A_-\). Thus the total-variation bound gives
\[
\begin{aligned}
Q_+^{n,\mathrm{rand}}(\mathcal A_+)
+Q_-^{n,\mathrm{rand}}(\mathcal A_-)
&\ge
Q_+^{n,\mathrm{rand}}(\mathcal A_+)
+Q_-^{n,\mathrm{rand}}(\mathcal A_+^c)\\
&\ge
1-\operatorname{TV}(Q_+^{n,\mathrm{rand}},Q_-^{n,\mathrm{rand}})
\ge\frac45,
\end{aligned}
\]
and consequently
\[
\max\left\{
Q_+^{n,\mathrm{rand}}(\mathcal A_+),
Q_-^{n,\mathrm{rand}}(\mathcal A_-)
\right\}\ge\frac25.
\]
Pointwise domination by the threshold times its indicator now implies
\[
\begin{aligned}
\max_{\pm}E_{Q_\pm^{n,\mathrm{rand}}}|T-\theta_\pm|
&\ge
\frac{\delta_{n,\sigma}}2
\max_{\pm}Q_\pm^{n,\mathrm{rand}}(\mathcal A_\pm)\\
&\ge\frac{\delta_{n,\sigma}}5.
\end{aligned}
\]
These are inequalities between nonnegative integrals, so they also apply to infinite risks. Since both witnesses belong to the model class, taking the infimum over estimators proves
\[
R_n(\sigma)\ge\frac{\delta_{n,\sigma}}5
=\frac{c_{\mathrm w}}5\rho_{n,\sigma}.
\]

Next, let \(C^\circ\) be any measurable connected-interval procedure honest over the model class, and define its coverage events by
\[
\mathcal G_\pm:=\{z:\theta_\pm\in C^\circ(z)\}.
\]
Honesty and total variation imply
\[
Q_{P_+}^n(\mathcal G_+)\ge1-\alpha,
\qquad
Q_{P_+}^n(\mathcal G_-)
\ge Q_{P_-}^n(\mathcal G_-)-\tau
\ge1-\alpha-\tau.
\]
The complementary-event union bound therefore yields
\[
Q_{P_+}^n(\mathcal G_+\cap\mathcal G_-)
\ge1-2\alpha-\tau.
\]
On this intersection, connectedness gives containment of the closed segment joining the targets:
\[
[\min\{\theta_+,\theta_-\},\max\{\theta_+,\theta_-\}]
\subseteq C^\circ(z).
\]
Its Lebesgue length is \(|\theta_+-\theta_-|\), so
\[
\operatorname{length}(C^\circ(z))
\ge
\delta_{n,\sigma}\mathbf1_{\mathcal G_+\cap\mathcal G_-}(z).
\]
Integrating gives
\[
\sup_{P\in\mathcal P_{K,\beta,\kappa,\sigma}}
E_{Q_P^n}\operatorname{length}(C^\circ)
\ge
E_{Q_{P_+}^n}\operatorname{length}(C^\circ)
\ge(1-2\alpha-\tau)\delta_{n,\sigma}.
\]
Taking the infimum over honest procedures proves the corresponding length bound. The argument includes \(\delta_{n,\sigma}=0\).

Define
\[
c_L:=c_{\mathrm w}\tau>0,
\qquad
n_L:=\max\{n_{\mathrm w},1\}.
\]
For \(n\ge n_L\), the inequalities \(\tau\le1/5\) and \(\tau\le1-2\alpha-\tau\), together with \(\delta_{n,\sigma}\ge0\), give
\[
c_L\rho_{n,\sigma}=\tau\delta_{n,\sigma}
\le\frac{\delta_{n,\sigma}}5\le R_n(\sigma),
\]
\[
c_L\rho_{n,\sigma}=\tau\delta_{n,\sigma}
\le(1-2\alpha-\tau)\delta_{n,\sigma}\le H_n(\sigma).
\]
We have obtained, uniformly over the path spaces and allowed error scales,
\[
c_L\rho_{n,\sigma}\le R_n(\sigma),
\qquad
c_L\rho_{n,\sigma}\le H_n(\sigma)
\qquad(n\ge n_L).
\]
% lean: CausalSmith.Stat.NoisydoseWeakdesignTransition.frontier_both_lower

\item \textbf{Comparison at the second threshold.}
Define the positive constant and logarithmic denominator by
\[
\lambda_*:=\log(e+1)>0,
\qquad
\Lambda_n:=\log(e+n b_n^d)\quad(n\ge1).
\]
Because \(L_n\to\infty\) and \(\log L_n/L_n\to0\), choose \(n_E\ge1\) such that, for every \(n\ge n_E\),
\[
L_n\ge4,
\qquad
\frac d2\frac{\log L_n}{L_n}<\frac12.
\]
Since \(b_n=L_n^{-1/2}\) and \(\log n=L_n-1\), logarithmic monotonicity gives
\[
\Lambda_n
\ge\log(n b_n^d)
=L_n-1-\frac d2\log L_n
>\frac{L_n}{2}-1
\ge\frac{L_n}{4}.
\]
Also \(0<b_n\le1\) and \(d\ge0\), whence
\[
e+n b_n^d\le e+n\le2en.
\]
Using \(\log2\le2\), we obtain
\[
\Lambda_n\le\log2+L_n\le4L_n.
\]
Taking square roots and multiplying by \(\sqrt{L_n}\) yields
\[
\frac{\sqrt{L_n}}2\le\sqrt{\Lambda_n}\le2\sqrt{L_n},
\qquad
\frac{L_n}{2}
\le\sqrt{L_n}\sqrt{\Lambda_n}
\le2L_n.
\]
The defining formulas consequently give
\[
F_n(b_n)=\frac1{\sqrt{L_n}\sqrt{\Lambda_n}},
\qquad
\frac1{2L_n}\le F_n(b_n)\le\frac2{L_n}.
\]
Furthermore, \(b_n^2L_n=1\), so
\[
P_n(b_n)=\frac{\lambda_*}{L_n}.
\]
Define
\[
C_E:=\max\left\{\frac2{\lambda_*},\,2\lambda_*\right\}>0.
\]
Then, for \(n\ge n_E\),
\[
F_n(b_n)
\le\frac2{L_n}
\le C_E\frac{\lambda_*}{L_n}
=C_EP_n(b_n),
\]
\[
P_n(b_n)
=\frac{\lambda_*}{L_n}
\le\frac{C_E}{2L_n}
\le C_EF_n(b_n).
\]
% lean: CausalSmith.Stat.NoisydoseWeakdesignTransition.secondElbow_comparison

\item \textbf{Uniform constants and both threshold comparisons.}
Set
\[
c:=c_L,
\qquad
C:=\max\left\{
C_R,C_H,\sqrt{\lambda_*},\frac1{\sqrt{\lambda_*}},
C_E,c_L+1
\right\},
\]
\[
n_0:=\max\{n_R,n_H,n_L,n_E\}.
\]
Thus \(0<c<C\) and \(n_0\ge1\). The second-threshold inequalities just proved remain valid with \(C\) in place of \(C_E\), because the comparison scales are nonnegative.

For \(n\ge n_0\), positivity of \(n\) and \(d\) gives the exact identities
\[
n a_n^d
=n(n^{-1/d})^d
=n n^{-1}=1,
\qquad
F_n(a_n)=\frac{a_n}{\sqrt{\lambda_*}}.
\]
The choices \(C\ge\sqrt{\lambda_*}\) and
\(C\ge1/\sqrt{\lambda_*}\) therefore imply
\[
D_n=a_n\le C F_n(a_n),
\qquad
F_n(a_n)\le C D_n.
\]
For every \(\sigma\in[0,1/4]\), nonnegativity of the branches in
\cref{obj:def:frontier-rate} gives \(\rho_{n,\sigma}\ge0\). Hence enlarging the upper-bound constants preserves their inequalities:
\[
c\rho_{n,\sigma}\le R_n(\sigma)
\le C_R\rho_{n,\sigma}\le C\rho_{n,\sigma},
\]
\[
c\rho_{n,\sigma}\le H_n(\sigma)
\le C_H\rho_{n,\sigma}\le C\rho_{n,\sigma}.
\]
Together with the second-threshold bounds, this establishes the first group of assertions with one \(C\) and one \(n_0\).
% lean: CausalSmith.Stat.NoisydoseWeakdesignTransition.uniform_frontier

\item \textbf{Comparison factors for transition windows.}
Fix a real \(\omega_{\mathrm{win}}\ge1\), and define
\[
M_{\mathrm F,\omega_{\mathrm{win}}}:=\omega_{\mathrm{win}}(1+d\log \omega_{\mathrm{win}}),
\qquad
M_{\mathrm P,\omega_{\mathrm{win}}}:=1+2\log \omega_{\mathrm{win}},
\]
\[
C_{\mathrm a}:=\max\left\{\sqrt{\lambda_*},
\frac1{\sqrt{\lambda_*}}\right\},
\qquad
C_K:=\max\{1,M_{\mathrm F,\omega_{\mathrm{win}}}C_{\mathrm a},
M_{\mathrm F,\omega_{\mathrm{win}}}C_EM_{\mathrm P,\omega_{\mathrm{win}}}\}.
\]
All these factors are positive, and \(C_K\ge1\).

We first derive the scale comparisons used throughout either window. For positive radii \(\varrho_1,\varrho_2\), a nonnegative exponent \(\gamma\), and a nonnegative weight \(\omega_{\log}\), suppose
\[
\varrho_1,\varrho_2>0,\qquad
\gamma,\omega_{\log}\ge0,\qquad
\varrho_1\le \omega_{\mathrm{win}}\varrho_2.
\]
Since \(\omega_{\mathrm{win}}^\gamma\ge1\) and
\(\varrho_1^\gamma\le \omega_{\mathrm{win}}^\gamma\varrho_2^\gamma\),
\[
e+\omega_{\log}\varrho_1^\gamma
\le \omega_{\mathrm{win}}^\gamma(e+\omega_{\log}\varrho_2^\gamma).
\]
Taking logarithms gives
\[
\log(e+\omega_{\log}\varrho_1^\gamma)
\le\gamma\log \omega_{\mathrm{win}}+
\log(e+\omega_{\log}\varrho_2^\gamma).
\]
The logarithm on the right is at least \(1\), while
\(\gamma\log \omega_{\mathrm{win}}\ge0\). Therefore,
\[
\log(e+\omega_{\log}\varrho_1^\gamma)
\le
(1+\gamma\log \omega_{\mathrm{win}})
\log(e+\omega_{\log}\varrho_2^\gamma).
\]
Since \(1+\gamma\log \omega_{\mathrm{win}}\ge1\), this also implies
\[
\log(e+\omega_{\log}\varrho_1^\gamma)
\le
(1+\gamma\log \omega_{\mathrm{win}})^2
\log(e+\omega_{\log}\varrho_2^\gamma),
\]
and taking nonnegative square roots yields
\[
\sqrt{\log(e+\omega_{\log}\varrho_1^\gamma)}
\le
(1+\gamma\log \omega_{\mathrm{win}})
\sqrt{\log(e+\omega_{\log}\varrho_2^\gamma)}.
\]

Now let \(n\ge1\), \(\varrho>0\), and
\(\sigma\in[\varrho/\omega_{\mathrm{win}},\omega_{\mathrm{win}}\varrho]\). Both
\(\sigma\le \omega_{\mathrm{win}}\varrho\) and \(\varrho\le \omega_{\mathrm{win}}\sigma\) hold.
Applying the square-root comparison in both directions with
\(\gamma=d\) and \(\omega_{\log}=n\) gives
\[
\begin{aligned}
F_n(\sigma)
&=
\frac{\sigma}{\varrho}
\frac{\sqrt{\log(e+n\varrho^d)}}
     {\sqrt{\log(e+n\sigma^d)}}F_n(\varrho)
\le M_{\mathrm F,\omega_{\mathrm{win}}}F_n(\varrho),\\
F_n(\varrho)
&=
\frac{\varrho}{\sigma}
\frac{\sqrt{\log(e+n\sigma^d)}}
     {\sqrt{\log(e+n\varrho^d)}}F_n(\sigma)
\le M_{\mathrm F,\omega_{\mathrm{win}}}F_n(\sigma).
\end{aligned}
\]
Applying the logarithmic comparison in both directions with
\(\gamma=2\) and \(\omega_{\log}=L_n>0\), and dividing by \(L_n\), gives
\[
P_n(\sigma)\le M_{\mathrm P,\omega_{\mathrm{win}}}P_n(\varrho),
\qquad
P_n(\varrho)\le M_{\mathrm P,\omega_{\mathrm{win}}}P_n(\sigma).
\]
Thus both scales vary by the stated fixed factors throughout any such radius window.
% lean: CausalSmith.Stat.NoisydoseWeakdesignTransition.fourierScale_window_comparison
% lean: CausalSmith.Stat.NoisydoseWeakdesignTransition.polynomialScale_window_comparison

\item \textbf{One constant and threshold for both windows.}
The positive exponent \(d\) implies
\[
a_n=n^{-1/d}\longrightarrow0,
\qquad
b_n=L_n^{-1/2}\longrightarrow0.
\]
Multiplication by the fixed \(\omega_{\mathrm{win}}\) preserves these limits. Choose
\(n_{\omega_{\mathrm{win}}}\in\mathbb N\) so that, for every \(n\ge n_{\omega_{\mathrm{win}}}\),
\[
\omega_{\mathrm{win}}a_n\le\frac14,\qquad \omega_{\mathrm{win}}b_n\le\frac14,
\]
and define
\[
n_K:=\max\{1,n_{\omega_{\mathrm{win}}},n_E\}.
\]
For \(n\ge n_K\), positivity of \(n,L_n,\omega_{\mathrm{win}}\) ensures
\[
0<\frac{a_n}{\omega_{\mathrm{win}}},\qquad
\omega_{\mathrm{win}}a_n\le\frac14,\qquad
0<\frac{b_n}{\omega_{\mathrm{win}}},\qquad
\omega_{\mathrm{win}}b_n\le\frac14.
\]

At the first threshold, the exact identity from the preceding argument gives
\[
D_n\le C_{\mathrm a}F_n(a_n),
\qquad
F_n(a_n)\le C_{\mathrm a}D_n.
\]
For \(\sigma\in[a_n/\omega_{\mathrm{win}},\omega_{\mathrm{win}}a_n]\), the window bounds therefore imply
\[
\begin{aligned}
D_n
&\le C_{\mathrm a}F_n(a_n)
\le C_{\mathrm a}M_{\mathrm F,\omega_{\mathrm{win}}}F_n(\sigma)
\le C_KF_n(\sigma),\\
F_n(\sigma)
&\le M_{\mathrm F,\omega_{\mathrm{win}}}F_n(a_n)
\le M_{\mathrm F,\omega_{\mathrm{win}}}C_{\mathrm a}D_n
\le C_KD_n.
\end{aligned}
\]
Since \(C_K>0\), these are precisely
\[
C_K^{-1}D_n\le F_n(\sigma)\le C_KD_n.
\]

For \(\sigma\in[b_n/\omega_{\mathrm{win}},\omega_{\mathrm{win}}b_n]\), combine both window comparisons with the second-threshold bounds:
\[
\begin{aligned}
P_n(\sigma)
&\le M_{\mathrm P,\omega_{\mathrm{win}}}P_n(b_n)
\le M_{\mathrm P,\omega_{\mathrm{win}}}C_EF_n(b_n)\\
&\le M_{\mathrm P,\omega_{\mathrm{win}}}C_EM_{\mathrm F,\omega_{\mathrm{win}}}F_n(\sigma)
\le C_KF_n(\sigma),\\
F_n(\sigma)
&\le M_{\mathrm F,\omega_{\mathrm{win}}}F_n(b_n)
\le M_{\mathrm F,\omega_{\mathrm{win}}}C_EP_n(b_n)\\
&\le M_{\mathrm F,\omega_{\mathrm{win}}}C_EM_{\mathrm P,\omega_{\mathrm{win}}}P_n(\sigma)
\le C_KP_n(\sigma).
\end{aligned}
\]
Dividing the first inequality by \(C_K\) establishes
\[
C_K^{-1}P_n(\sigma)\le F_n(\sigma)\le C_KP_n(\sigma).
\]
The choices of \(C_K,n_K\) precede the choice of \(\sigma\), and the same choices apply to both entire windows.
% lean: CausalSmith.Stat.NoisydoseWeakdesignTransition.frontier_transition_windows

\item \textbf{Fixed positive noise and the frontier rate.}
Fix \(\sigma\in(0,1/4]\). Since both thresholds tend to zero, eventually
\[
a_n<\sigma,\qquad b_n<\sigma.
\]
The two branch tests in \cref{obj:def:frontier-rate} then select
\[
b_{n,\sigma}=P_n(\sigma)
=\frac{\log(e+\sigma^2L_n)}{L_n}.
\]

For natural \(n\) with \(\log n\ge1\), define
\[
\zeta_n:=\frac{\log\log n}{\log n}.
\]
Both \(\log n\) and \(\log\log n\) tend to infinity. Hence eventually
\[
\log n\ge1,\qquad
\log\log n\ge0,\qquad
\log\log n\ge-2\log(\sigma^2),
\qquad
\log\log n\ge\log(e+2\sigma^2).
\]
Also \(L_n=1+\log n>0\). Monotonicity of the logarithm gives the lower numerator bound
\[
\begin{aligned}
\log(e+\sigma^2L_n)
&\ge\log(\sigma^2\log n)\\
&=\log(\sigma^2)+\log\log n
\ge\frac12\log\log n.
\end{aligned}
\]
For the upper bound, \(\log n\ge1\) implies
\[
e+\sigma^2L_n
=e+\sigma^2(1+\log n)
\le(e+2\sigma^2)\log n.
\]
Thus
\[
\log(e+\sigma^2L_n)
\le\log(e+2\sigma^2)+\log\log n
\le2\log\log n.
\]
We also have
\[
0\le\zeta_n\le\log\log n.
\]
Consequently,
\[
\frac14\zeta_nL_n
=\frac14\log\log n+\frac14\zeta_n
\le\frac12\log\log n
\le\log(e+\sigma^2L_n),
\]
and
\[
\log(e+\sigma^2L_n)
\le2\log\log n
\le2\log\log n+2\zeta_n
=2\zeta_nL_n.
\]
Dividing by \(L_n>0\) proves
\[
0\le\zeta_n,\qquad
\frac14\zeta_n\le P_n(\sigma)\le2\zeta_n.
\]
Since \(\beta\ge0\), taking the \(\beta\)-th power preserves these comparisons. Together with the selected branch, this gives
\[
0\le\zeta_n^\beta,\qquad
4^{-\beta}\zeta_n^\beta
\le\rho_{n,\sigma}
\le2^\beta\zeta_n^\beta.
\]
Choose \(n_F\in\mathbb N\) so that all these assertions hold for every \(n\ge n_F\).
% lean: CausalSmith.Stat.NoisydoseWeakdesignTransition.fixedNoise_frontierRate_bounds

\item \textbf{Fixed-noise risk and length constants.}
Define
\[
c_\sigma:=c_L4^{-\beta},
\qquad
C_\sigma:=
\max\left\{\max\{C_R,C_H\}2^\beta,c_\sigma\right\}+1,
\]
\[
n_\sigma:=\max\{n_F,n_L,n_R,n_H\}.
\]
Then \(0<c_\sigma<C_\sigma\). For every public path space and every \(n\ge n_\sigma\), the lower-rate comparison and the common lower bound yield
\[
c_\sigma\zeta_n^\beta
\le c_L\rho_{n,\sigma}\le R_n(\sigma),
\qquad
c_\sigma\zeta_n^\beta
\le c_L\rho_{n,\sigma}\le H_n(\sigma).
\]
The upper-rate comparison and the two upper bounds yield
\[
R_n(\sigma)
\le C_R\rho_{n,\sigma}
\le C_R2^\beta\zeta_n^\beta
\le C_\sigma\zeta_n^\beta,
\]
\[
H_n(\sigma)
\le C_H\rho_{n,\sigma}
\le C_H2^\beta\zeta_n^\beta
\le C_\sigma\zeta_n^\beta.
\]
Substituting the definition of \(\zeta_n\) proves the fixed-positive-noise assertions. Every constant and sample threshold was chosen independently of the public path space, as required.
% lean: CausalSmith.Stat.NoisydoseWeakdesignTransition.uniform_frontier
\end{enumerate}
\end{proof}
